\chapter{Supercharges, Closure, and Superspace}

\lectureoverview{Supersymmetry is difficult not because its rules are numerous, but because the rules join ideas that ordinary physics keeps separate. We begin with Grassmann variables, whose odd elements anticommute and whose functions terminate after finitely many terms. Their derivative algebra then provides the model for odd symmetry generators. A bosonic oscillator and a fermionic oscillator give the simplest supercharges: their anticommutator is the Hamiltonian, so exchanging bosons and fermions is inseparable from time translation. The same algebra can be represented geometrically by translations in Grassmann directions of an enlarged spacetime called superspace.}

\section{Why We Must Back Up Before Going Forward}

Supersymmetry is among the most abstract constructions in physics. It resembles rotations and translations because it is a symmetry, but it does not fit inside ordinary spacetime without enlarging what we mean by a coordinate. Its natural parameters are Grassmann quantities: objects that can be added and multiplied, yet anticommute rather than commute.

That is why a direct assault on the final algebra would be unhelpful. We shall first make the odd arithmetic usable, then recall how ordinary continuous symmetries are organized, and only then combine the two ideas. The route is longer, but each step prepares the next.

\section{Grassmann Algebra in Working Form}

\subsection{Odd and even elements}

Let \(\epsilon_i\) denote independent Grassmann generators. They obey
\begin{equation}
\epsilon_i\epsilon_j=-\epsilon_j\epsilon_i,
\qquad
\{\epsilon_i,\epsilon_j\}=0.
\end{equation}
Setting \(i=j\) gives
\begin{equation}
2\epsilon_i^2=0,
\qquad
\epsilon_i^2=0,
\end{equation}
over the real or complex numbers. The generator is not zero; only its square is zero.

Each homogeneous element has a parity. A product containing an odd number of Grassmann generators is odd, while a product containing an even number is even:
\begin{align}
\text{odd}\times\text{odd}&=\text{even},\\
\text{odd}\times\text{even}&=\text{odd},\\
\text{even}\times\text{even}&=\text{even}.
\end{align}
Even elements commute with homogeneous elements. Two odd elements anticommute. Thus
\begin{equation}
n\epsilon_i=\epsilon_i n,
\qquad
\epsilon_i\epsilon_j=-\epsilon_j\epsilon_i,
\end{equation}
when \(n\) is even.

An even element need not be an ordinary real or complex number. The product
\(\epsilon_1\epsilon_2\), for example, is even and nilpotent. Ordinary numbers form the scalar part of the algebra, while even products provide further commuting elements.

\begin{classroomqa}
\audiencequestion{Is there a distinction between the Grassmann algebra and the Grassmann numbers, or are they the same thing? \lecturetimestamp{00:06:45}}
\lecturerresponse{The elements are the Grassmann numbers; the algebra is the collection of those elements together with their addition and multiplication rules. One may therefore speak loosely of the same object, but it is useful to distinguish the things being manipulated from the rules by which they are manipulated.}
\end{classroomqa}

\subsection{Functions terminate}

A function of one Grassmann variable has only two possible powers:
\begin{equation}
f(\epsilon)=a+b\epsilon.
\end{equation}
There is no \(\epsilon^2\) term. With two independent generators,
\begin{equation}
f(\epsilon_1,\epsilon_2)
=a+b_1\epsilon_1+b_2\epsilon_2+c\epsilon_1\epsilon_2.
\end{equation}
The expansion stops after both generators have appeared. For \(n\) independent generators, the ordered monomials form a \(2^n\)-dimensional basis:
\begin{equation}
1,\quad
\epsilon_i,\quad
\epsilon_i\epsilon_j,\quad
\ldots,\quad
\epsilon_1\epsilon_2\cdots\epsilon_n.
\end{equation}
\editorialnote{The finite \(2^n\)-dimensional statement is the general pattern behind the one- and two-generator examples. An infinite Grassmann algebra can also be defined, but every polynomial used here involves only finitely many generators.}

When two real generators are combined into a complex pair \(\theta,\bar\theta\), a general function can be written, with a chosen ordering, as
\begin{equation}
f(\theta,\bar\theta)
=a+\bar b\,\theta+\bar\theta\,b+c\,\bar\theta\theta.
\end{equation}
The bar indicates the complex-pair notation used here. Reversing the final order changes a sign:
\begin{equation}
\bar\theta\theta=-\theta\bar\theta.
\end{equation}

It is usually convenient to assign the whole function a definite parity. If \(f\) is even, then \(a\) and \(c\) are even while \(b\) and \(\bar b\) are odd. If \(f\) is odd, those assignments reverse. The coefficients may themselves belong to a larger Grassmann algebra, provided they do not depend on the particular variables with respect to which we are expanding.

\begin{classroomqa}
\audiencequestion{Can a Grassmann variable be inserted into trigonometric or other analytic functions? \lecturetimestamp{00:11:06}}
\lecturerresponse{Yes, but the power series terminates immediately. For one odd generator,
\[
\sin\epsilon=\epsilon,\qquad
\cos\epsilon=1,\qquad
e^\epsilon=1+\epsilon.
\]
The familiar names survive, but there are too few powers for the functions to develop their usual analytic complexity.}
\end{classroomqa}

\section{Differentiation Is Sign Bookkeeping}

We use a left Grassmann derivative. Its defining local rules are
\begin{equation}
\partial_\theta a=0,
\qquad
\partial_\theta(\theta B)=B,
\end{equation}
when \(a\) and \(B\) are independent of \(\theta\). Before differentiating, move the relevant \(\theta\) to the left. Every passage through another odd factor changes the sign.

For example, if \(b\) is odd,
\begin{equation}
b\theta=-\theta b,
\qquad
\partial_\theta(b\theta)=-b.
\end{equation}
Nothing mysterious has happened to the derivative. The minus sign arose before differentiation, when the odd factors were reordered.

\subsection{The two-variable calculation}

Return to
\begin{equation}
f(\theta,\bar\theta)
=a+\bar b\,\theta+\bar\theta\,b+c\,\bar\theta\theta.
\end{equation}
Suppose first that \(f\) is even. Then \(\bar b\) is odd and \(c\) is even. Reordering the terms that contain \(\theta\),
\begin{align}
\bar b\,\theta&=-\theta\bar b,\\
c\,\bar\theta\theta&=-\theta c\,\bar\theta.
\end{align}
Therefore
\begin{equation}
\partial_\theta f=-\bar b-c\bar\theta.
\end{equation}

If \(f\) is odd, then \(\bar b\) is even and \(c\) is odd. The first term crosses no odd coefficient, while the last \(\theta\) crosses two odd factors, \(c\) and \(\bar\theta\). The two signs cancel:
\begin{equation}
\partial_\theta f=\bar b+c\bar\theta.
\end{equation}

\begin{classroomqa}
\audiencequestion{Why do both surviving terms carry minus signs in the even example? \lecturetimestamp{00:16:49}}
\lecturerresponse{The first sign comes from moving \(\theta\) past the odd coefficient \(\bar b\). In the last term, \(c\) is even, but \(\theta\) must pass the odd variable \(\bar\theta\). Each reordering crosses one odd factor, so each produces one minus sign.}
\end{classroomqa}

The answer depends on the ordering convention, but the mathematics does not. Had we begun with \(c\theta\bar\theta\) rather than \(c\bar\theta\theta\), the displayed coefficient would differ by a sign and the derivative formula would change correspondingly.

\begin{classroomqa}
\audiencequestion{Could the last term have been written with \(\theta\) and \(\bar\theta\) reversed, and does the resulting sign have a deeper significance? \lecturetimestamp{00:18:50}}
\lecturerresponse{It can be written in either order, but exchanging the two odd variables changes the sign. One may absorb that sign into the coefficient \(c\). At this stage the effect is bookkeeping: order matters, parity does not change, and a sign change never turns an odd object into an even one. The physical significance comes when the same bookkeeping enforces fermionic statistics and supersymmetry.}
\end{classroomqa}

\begin{classroomqa}
\audiencequestion{Are the coefficients \(a,b,c\) simply drawn from the complex-number field? \lecturetimestamp{00:20:47}}
\lecturerresponse{They may be ordinary complex numbers when their required parity is even, but they may also contain generators other than the selected \(\theta,\bar\theta\). The algebra of interest normally contains all complex scalars together with a chosen set of odd generators and their products.}
\end{classroomqa}

\begin{classroomqa}
\audiencequestion{The derivatives of the even and odd functions look like opposites. Is that a general symmetry? \lecturetimestamp{00:22:31}}
\lecturerresponse{Not by itself. The signs follow mechanically from the parity of each coefficient and the number of odd factors crossed. Sometimes two answers happen to be opposite, but the reliable rule is simpler: every exchange of two odd objects contributes a minus sign.}
\end{classroomqa}

\subsection{Coordinate and derivative as operators}

There is an instructive ordinary analogue. Let multiplication by \(x\) and differentiation by \(x\) act on an arbitrary function \(g(x)\). Then
\begin{align}
\partial_x\!\left(xg\right)&=g+xg',\\
x\partial_xg&=xg',
\end{align}
so
\begin{equation}
[\partial_x,x]g=g.
\end{equation}
Because this holds for every \(g\),
\begin{equation}
[\partial_x,x]=1.
\end{equation}
With \(p=-i\partial_x\), this is the familiar position-momentum algebra, in units with \(\hbar=1\).

The Grassmann analogue uses an anticommutator. Take
\begin{equation}
f(\theta)=a+\theta b,
\end{equation}
with \(a\) even and \(b\) odd, so that \(f\) is even. First,
\begin{equation}
\theta\partial_\theta f=\theta b.
\end{equation}
In the opposite order,
\begin{align}
\partial_\theta(\theta f)
&=\partial_\theta(\theta a+\theta^2b)\\
&=a.
\end{align}
Adding the two orders gives the original function:
\begin{equation}
\left(\partial_\theta\theta+\theta\partial_\theta\right)f
=a+\theta b=f.
\end{equation}
Hence
\begin{equation}
\boxed{\{\partial_\theta,\theta\}=1.}
\end{equation}
The derivative \(\partial_\theta\) is itself odd. That is why its canonical relation with the odd coordinate is an anticommutator.

\begin{classroomqa}
\audiencequestion{May one use a function containing both even and odd parts rather than assigning a definite parity? \lecturetimestamp{00:31:05}}
\lecturerresponse{Such inhomogeneous elements can be formed, but calculations are normally organized into homogeneous even and odd objects. A definite parity makes every sign in a graded product unambiguous; a mixed object is handled by splitting it into its even and odd parts.}
\end{classroomqa}

\begin{classroomqa}
\audiencequestion{Was the test function \(g(x)\) in the ordinary commutator calculation supposed to be odd? \lecturetimestamp{00:31:45}}
\lecturerresponse{No. That was an ordinary function of an ordinary commuting coordinate. It was introduced only to show how the operator identity \([\partial_x,x]=1\) is proved before repeating the argument with the Grassmann pair \(\partial_\theta,\theta\).}
\end{classroomqa}

\begin{classroomqa}
\audiencequestion{Is there a limit to the number of independent Grassmann generators? \lecturetimestamp{00:32:16}}
\lecturerresponse{There is no algebraic upper limit. For \(n\) generators the basis has \(2^n\) monomials, so the bookkeeping grows rapidly, but the construction itself works for any finite \(n\).}
\end{classroomqa}

\subsection{Integration}

The integral over one Grassmann variable is fixed by
\begin{equation}
\int d\theta\,1=0,
\qquad
\int d\theta\,\theta=1.
\end{equation}
Consequently, for \(f(\theta)=a+\theta b\),
\begin{equation}
\int d\theta\,f(\theta)=b=\partial_\theta f(\theta).
\end{equation}
Thus integration and differentiation perform the same component extraction. This is the Berezin integral. Its normalization is the unique simple choice compatible with translation invariance and Grassmann integration by parts. \editorialnote{The component rules make precise the lecture's statement that Grassmann integration and differentiation are the same operation.}

\section{Ordinary Symmetry as the Baseline}

Consider a continuous unitary group with Hermitian generators \(G_i\). Near the identity,
\begin{equation}
U(\alpha)=\mathbf 1+i\alpha^iG_i+O(\alpha^2),
\end{equation}
where the \(\alpha^i\) are ordinary commuting parameters. The infinitesimal multiplication law is encoded by the Lie algebra
\begin{equation}
[G_i,G_j]=i f_{ij}{}^{k}G_k.
\end{equation}
The constants \(f_{ij}{}^k\) are the structure constants. Closure means that commuting any two generators returns a linear combination of generators already in the algebra.

For rotations,
\begin{equation}
[L_i,L_j]=i\varepsilon_{ijk}L_k.
\end{equation}
The three Hermitian components may be combined into the non-Hermitian ladder operators
\begin{equation}
L_\pm=L_x\pm iL_y,
\end{equation}
which raise or lower the magnetic quantum number \(m\) by one. Complex combinations are useful without changing the Hermiticity of the underlying physical generators.

Only the first two indices of \(f_{ij}{}^k\) must be antisymmetric in a general basis. For a compact algebra with an invariant orthonormal metric, one can lower the final index and choose conventions in which \(f_{ijk}\) is totally antisymmetric, as it is for rotations. \editorialnote{This qualification separates the universally antisymmetric commutator from the stronger property available in the familiar rotation basis.}

For \(G_i\) to generate an actual symmetry of a time-independent system,
\begin{equation}
[H,G_i]=0.
\end{equation}
Equivalently, \(G_i\) maps every energy eigenspace into itself. With no explicit time dependence, the Heisenberg equation gives
\begin{equation}
\frac{dG_i}{dt}=i[H,G_i]=0,
\end{equation}
so the generator is a conserved charge.

Ordinary observables are even. A fermion field is odd, but an isolated odd field is not itself a measurable observable; physical operators contain an even total number of fermionic factors. Supersymmetry changes the parity of the generators, not the parity required of observable results.

The clean notation for all cases is the graded commutator
\begin{equation}
[X,Y]_{\mathrm{gr}}
=XY-(-1)^{|X||Y|}YX.
\end{equation}
It is an ordinary commutator if either operator is even and an anticommutator if both are odd.

\section{The Smallest Supersymmetric Oscillator}

\subsection{One bosonic mode and one fermionic mode}

Take a single bosonic mode with creation and annihilation operators \(a_+\) and \(a_-\):
\begin{equation}
[a_-,a_+]=1,
\qquad
[a_+,a_+]=[a_-,a_-]=0.
\end{equation}
Its number operator is
\begin{equation}
N_B=a_+a_-,
\end{equation}
and its occupation number can be \(0,1,2,\ldots\).

For a single fermionic mode, introduce \(c_+\) and \(c_-\):
\begin{equation}
\{c_-,c_+\}=1,
\qquad
\{c_+,c_+\}=\{c_-,c_-\}=0.
\end{equation}
Therefore
\begin{equation}
c_+^2=c_-^2=0,
\qquad
N_F=c_+c_-,
\end{equation}
and the occupation number is only \(0\) or \(1\). The bosonic operators are even, the fermionic operators are odd, and every \(a_\pm\) commutes with every \(c_\pm\).

\editorialnote{The nilpotence restricts the fermionic mode to one occupied state. The bosonic mode has unrestricted occupation; a brief contrary phrase in the classroom discussion is resolved by the oscillator algebra itself.}

Assume the boson and fermion have the same rest energy \(m\), and define
\begin{equation}
Q=\sqrt m\,a_+c_-,
\qquad
\bar Q=Q^\dagger=\sqrt m\,a_-c_+.
\end{equation}
The operator \(Q\) removes one fermion and creates one boson. Its conjugate removes one boson and creates one fermion. Each contains one fermionic operator and is therefore odd.

On number states, their action is explicit:
\begin{align}
Q\ket{n,1}
&=\sqrt{m(n+1)}\,\ket{n+1,0},
&
Q\ket{n,0}&=0,\\
\bar Q\ket{n,0}
&=\sqrt{mn}\,\ket{n-1,1},
&
\bar Q\ket{n,1}&=0.
\end{align}
\editorialnote{The number-state formulas are the direct normalized form of the particle-replacement action described in the lecture.}

\subsection{The anticommutator closes on energy}

The graded analogue of finding a Lie algebra is to compute
\begin{align}
\{Q,\bar Q\}
&=Q\bar Q+\bar Q Q\\
&=m\left(a_+c_-a_-c_+
        +a_-c_+a_+c_-\right).
\end{align}
Since the \(a\)'s are even, they pass through the \(c\)'s without a sign:
\begin{equation}
\{Q,\bar Q\}
=m\left(a_+a_-c_-c_+
        +a_-a_+c_+c_-\right).
\end{equation}
Now use
\begin{equation}
c_-c_+=1-c_+c_-,
\qquad
a_-a_+=a_+a_-+1.
\end{equation}
Then
\begin{align}
\{Q,\bar Q\}
&=m\left[
N_B(1-N_F)+(N_B+1)N_F
\right]\\
&=m(N_B+N_F).
\end{align}
The temporary extra factor of two that arose during the live rearrangement disappears: one contribution supplies \(N_B\), the other supplies \(N_F\).

For particles at rest with equal masses, the Hamiltonian is
\begin{equation}
H=m(N_B+N_F).
\end{equation}
Thus
\begin{equation}
\boxed{\{Q,\bar Q\}=H.}
\end{equation}

\begin{figure}[H]
\centering
\includegraphics[width=0.62\textwidth]{lecture_05_figure_02.png}
\lectureframe{00:55:46.300}
\caption{The central closure relation isolated on the board. The braces and \(Q,\bar Q,H\) are visible; the comma is supplied by standard anticommutator notation.}
\end{figure}

The additive zero of energy has been chosen so that the empty state has \(H=0\). In the supersymmetric oscillator this is the natural convention: the bosonic and fermionic zero-point contributions cancel. \editorialnote{The lecture works directly with total rest energy and does not display separate zero-point terms; the stated convention makes that choice explicit.}

The Hamiltonian is even. It is also a symmetry generator: it generates translations in time. The odd generators therefore fail to close among themselves unless the even time-translation generator is included.

\begin{classroomqa}
\audiencequestion{Which member of the relation is even? Is it the Hamiltonian? \lecturetimestamp{00:58:24}}
\lecturerresponse{Yes. Energy is an ordinary observable, so \(H\) is even. The supercharge \(Q\) is odd, and the product of two odd quantities is even; accordingly, \(\{Q,\bar Q\}\) can equal \(H\). An odd and an even generator are related by a commutator, not an anticommutator.}
\end{classroomqa}

\subsection{Completing the algebra}

The same-parity relations are immediate:
\begin{equation}
\{Q,Q\}=2Q^2=0,
\qquad
\{\bar Q,\bar Q\}=2\bar Q^2=0,
\end{equation}
because \(Q^2\) contains \(c_-^2\) and \(\bar Q^2\) contains \(c_+^2\).

What remains is the relation with \(H\). Let
\begin{equation}
H\ket{E}=E\ket{E}.
\end{equation}
If the boson and fermion masses are equal, \(Q\) replaces one by the other without changing the energy. Hence \(Q\ket{E}\) is either zero or another eigenstate with the same eigenvalue:
\begin{equation}
H Q\ket{E}=E Q\ket{E}.
\end{equation}
On the other hand,
\begin{equation}
QH\ket{E}=Q(E\ket{E})=E Q\ket{E}.
\end{equation}
Subtracting,
\begin{equation}
[H,Q]\ket{E}=0.
\end{equation}
If the energy eigenstates form a complete basis, this holds as an operator identity:
\begin{equation}
[H,Q]=0,
\qquad
[H,\bar Q]=0.
\end{equation}

\begin{classroomqa}
\audiencequestion{What happens if the boson and fermion do not have the same mass? \lecturetimestamp{01:01:16}}
\lecturerresponse{Then replacing one by the other changes the energy, so the proposed supercharge does not commute with the Hamiltonian and the simple supersymmetry is lost. For
\[
H=m_BN_B+m_FN_F,
\]
the explicit result is
\[
[H,Q]=(m_B-m_F)Q.
\]
It vanishes precisely when \(m_B=m_F\). \editorialnote{The unequal-mass commutator is the direct operator form of the energy-change argument and makes the equal-mass condition explicit.}}
\end{classroomqa}

\begin{classroomqa}
\audiencequestion{Could the energy-eigenstate argument be repeated more slowly? \lecturetimestamp{01:01:37}}
\lecturerresponse{Begin with \(H\ket{E}=E\ket{E}\). Equal masses imply that \(Q\ket{E}\) has the same definite energy \(E\), so \(HQ\ket{E}=EQ\ket{E}\). The other order gives \(QH\ket{E}=QE\ket{E}=EQ\ket{E}\), because \(E\) is an ordinary scalar. Their difference vanishes on every energy eigenstate. Completeness then promotes the statement from those states to the operator equation \([H,Q]=0\).}
\end{classroomqa}

The complete toy algebra is therefore
\begin{align}
\{Q,\bar Q\}&=H,\\
\{Q,Q\}&=\{\bar Q,\bar Q\}=0,\\
[H,Q]&=[H,\bar Q]=0.
\end{align}
The supercharges are conserved odd charges. More surprisingly, their algebra mixes an operation that changes particle type with the generator of spacetime translation.

\section{Why the Hamiltonian Means Time Translation}

In this stripped-down model there is no spatial coordinate, but there is time. The Schrödinger equation reads
\begin{equation}
i\frac{\partial}{\partial t}\ket{\psi(t)}
=H\ket{\psi(t)}.
\end{equation}

\begin{figure}[H]
\centering
\includegraphics[width=0.56\textwidth]{lecture_05_figure_03.png}
\lectureframe{01:07:26.300}
\caption{The time derivative being written beside \(\ket{\psi}\). The hand obscures part of the denominator; the complete Schrödinger equation is given above.}
\end{figure}

For a finite shift \(\delta\),
\begin{equation}
\ket{\psi(t+\delta)}
=e^{-iH\delta}\ket{\psi(t)}.
\end{equation}
Expanding to first order,
\begin{equation}
\ket{\psi(t+\delta)}
=\left(\mathbf 1-i\delta H\right)\ket{\psi(t)}
+O(\delta^2).
\end{equation}
Thus \(H\) is the infinitesimal generator that updates a state in time.

Internal symmetries such as isospin and color also change particle labels. Isospin can mix proton and neutron states; color rotations mix red, green, and blue quarks. Yet no composition of those internal transformations translates a system in space or time. Supersymmetry is different:
\begin{equation}
\begin{gathered}
\text{odd particle-changing transformation}\\
\mathrel{\stackrel{\{\ ,\ \}}{\Longrightarrow}}\\
\text{time translation}
\end{gathered}
\end{equation}
That is the clue that the odd generators may themselves be translations, but along coordinates unlike the usual ones.

\section{Superspace}

\subsection{A function of time and odd coordinates}

Introduce a function
\begin{equation}
\Psi=\Psi(\theta,\bar\theta,t).
\end{equation}
It may be regarded as a wavefunction in the toy model or, more generally, as a superfield. Because the odd expansion terminates,
\begin{equation}
\Psi(\theta,\bar\theta,t)
=A(t)+\bar\chi(t)\theta
+\bar\theta\chi(t)+F(t)\bar\theta\theta.
\end{equation}
The single object packages even and odd component functions. \editorialnote{This component expansion applies the Grassmann polynomial developed earlier; it is the standard explicit content of a function on the toy superspace.}

An ordinary time translation leaves \(\theta,\bar\theta\) fixed and sends
\begin{equation}
t\longmapsto t+\delta.
\end{equation}
To first order,
\begin{equation}
\Psi(\theta,\bar\theta,t+\delta)
=\Psi(\theta,\bar\theta,t)
+\delta\,\partial_t\Psi(\theta,\bar\theta,t).
\end{equation}

\begin{figure}[H]
\centering
\includegraphics[width=0.72\textwidth]{lecture_05_figure_04.png}
\lectureframe{01:13:15.060}
\caption{The board places \(t\to t+\delta\) above the first-order change in \(\Psi\), with the derivative term indicated directly.}
\end{figure}

\subsection{Translation in an odd direction}

Now invent a transformation that shifts an odd coordinate:
\begin{equation}
\theta\longmapsto\theta+\epsilon,
\end{equation}
where \(\epsilon\) is another odd Grassmann parameter. At the same time, let time acquire an even nilpotent displacement,
\begin{equation}
t\longmapsto t+i\bar\theta\epsilon.
\end{equation}
Both \(\bar\theta\) and \(\epsilon\) are odd, so their product is even and can consistently be added to the even coordinate \(t\).

With the parameter moved to the left,
\begin{equation}
i\bar\theta\epsilon=-i\epsilon\bar\theta.
\end{equation}
The induced variation is therefore
\begin{align}
\delta_\epsilon\Psi
&=\epsilon\,\partial_\theta\Psi
-i\epsilon\bar\theta\,\partial_t\Psi\\
&=\epsilon q_0\Psi,
\end{align}
where
\begin{equation}
q_0=\partial_\theta-i\bar\theta\partial_t.
\end{equation}
This odd differential operator is the generator of the combined shift. A compatible conjugate generator, using left derivatives, is
\begin{equation}
\bar q_0=-\partial_{\bar\theta}
+i\theta\partial_t.
\end{equation}

The exact signs depend on whether one uses left or right Grassmann derivatives and on how conjugation reverses odd factors. What must not depend on convention is the closure on time translation. With the definitions above,
\begin{align}
\{q_0,\bar q_0\}
&=i\{\partial_\theta,\theta\}\partial_t
+i\{\bar\theta,\partial_{\bar\theta}\}\partial_t\\
&=2i\partial_t.
\end{align}
The derivative-derivative term vanishes, as does the term quadratic in the odd coordinates. The two mixed terms each contribute \(i\partial_t\).

Since \(H=i\partial_t\) on Schrödinger-picture states,
\begin{equation}
\{q_0,\bar q_0\}=2H.
\end{equation}
If we normalize
\begin{equation}
q=\frac{q_0}{\sqrt2},
\qquad
\bar q=\frac{\bar q_0}{\sqrt2},
\end{equation}
then
\begin{equation}
\boxed{\{q,\bar q\}=H.}
\end{equation}
\editorialnote{The lecture deliberately works through uncertain signs and then notices a factor of two. The displayed left-derivative convention preserves that classroom calculation while fixing one self-consistent normalization. Other standard conventions move signs between the coordinate shifts and the generators.}

\begin{classroomqa}
\audiencequestion{Where does the final factor of two in the superspace anticommutator come from? \lecturetimestamp{01:24:21}}
\lecturerresponse{There are two equal mixed contributions. One comes from \(\{\partial_\theta,\theta\}=1\), and the other from \(\{\partial_{\bar\theta},\bar\theta\}=1\). Each multiplies the time derivative. The unnormalized generators therefore give \(2i\partial_t\); a factor \(1/\sqrt2\) in each supercharge restores the earlier normalization.}
\end{classroomqa}

\begin{classroomqa}
\audiencequestion{Did the construction use the assumption that \(\epsilon\) is small? \lecturetimestamp{01:24:51}}
\lecturerresponse{It used nilpotence rather than an ordinary numerical magnitude. Since \(\epsilon^2=0\), every expansion in a single \(\epsilon\) stops after first order. Calling it infinitesimal or small is a useful analogy, but the truncation is exact.}
\end{classroomqa}

\editorialnote{Grassmann parameters have no ordinary positive magnitude. Their role as infinitesimal parameters comes from nilpotence and the exact termination of the expansion.}

\subsection{The geometric meaning}

The algebra now has two equivalent realizations. In the oscillator picture, \(Q\) and \(\bar Q\) replace fermions by bosons and bosons by fermions. In the geometric picture, \(q\) and \(\bar q\) translate a function through odd directions while inducing an ordinary time shift. In both pictures,
\begin{equation}
\begin{gathered}
\{\text{supercharge},\text{conjugate supercharge}\}\\
=\text{time-translation generator}.
\end{gathered}
\end{equation}

The generalized coordinate space containing \(t,\theta,\bar\theta\) is superspace. In a relativistic field theory it also contains the ordinary spatial coordinates, and the supercharges close on the full spacetime momentum \(P_\mu\), not only on \(H=P_0\). Curving this structure leads toward supergravity.

In one sense Grassmann variables are an efficient bookkeeping device for bosons and fermions. In another, the closure relation says something much deeper: the distinction between particle type and spacetime geometry is no longer absolute. Moving twice in odd directions produces motion in an ordinary direction. Whether nature realizes that structure experimentally is separate from the mathematical fact that the structure is coherent, economical, and remarkably rigid.
